\begin{ouverture}[probabilites]{11}{Probabilités : événements et deux épreuves}{Partie 3 · Données et hasard}
{Univers · Contraire · Réunion · Intersection · Fréquences}
\objectif{Décrire une expérience aléatoire avec le vocabulaire et les notations des ensembles.}
\objectif{Calculer la probabilité d’un événement et de son contraire.}
\objectif{Étudier une expérience à deux épreuves et observer la fluctuation des fréquences.}
\prerequis{Vocabulaire des probabilités et équiprobabilité (Cinquième), fractions.}
\end{ouverture}
\pagestyle{otMain}\essentiel\fontsize{11}{15}\selectfont
\phantomsection\label{cours-11-1}
\essSection{1}{Issues, événements et ensembles}
\begin{definition}Une expérience aléatoire a plusieurs résultats possibles, les issues, sans que l’on puisse prévoir lequel se produira. L’ensemble de toutes les issues est l’univers : pour un dé à six faces, $\{1\,;2\,;3\,;4\,;5\,;6\}$. Un événement est un ensemble d’issues, par exemple $A=\{3\,;6\}$, « obtenir un multiple de 3 ».\end{definition}
\begin{definition}[Opérations sur les événements]Le contraire de $A$, noté $\overline{A}$, contient les issues qui ne sont pas dans $A$ : c’est son complémentaire. L’intersection $A\cap B$, « $A$ et $B$ », contient les issues communes. La réunion $A\cup B$, « $A$ ou $B$ », contient les issues d’au moins un des deux événements. L’ensemble vide, noté $\varnothing$, est l’événement impossible.\end{definition}
\begin{propriete}La probabilité $P(A)$ d’un événement est un nombre compris entre 0 et 1 : c’est la somme des probabilités des issues qui le composent. Si les issues sont équiprobables, $P(A)$ est le nombre d’issues de $A$ divisé par le nombre total d’issues. Enfin $P(\overline A)=1-P(A)$.\end{propriete}
\begin{exemple}Dé équilibré, $A=\{3\,;6\}$ et $B=\{2\,;4\,;6\}$ : $A\cap B=\{6\}$, $A\cup B=\{2\,;3\,;4\,;6\}$ et $\overline A=\{1\,;2\,;4\,;5\}$. Donc $P(A\cup B)=\frac46=\frac23$ et $P(\overline A)=1-\frac26=\frac23$.\end{exemple}
\clearpage
\phantomsection\label{cours-11-2}
\essSection{2}{Deux épreuves et fréquences}
\begin{savoirfaire}[Lister les issues]Pour une expérience à deux épreuves, on liste les couples d’issues avec un arbre ou un tableau à double entrée. Une pièce puis un dé donnent $2\times6=12$ issues, de $(P\,;1)$ à $(F\,;6)$ ; avec une pièce et un dé équilibrés, on choisit l’équiprobabilité de ces 12 issues.\end{savoirfaire}
\label{fig-cours-11-2-cours-arbre-pieces}\figureGeo{cours-arbre-pieces}

\begin{exemple}Deux pièces équilibrées que l’on distingue : les quatre issues $PP$, $PF$, $FP$, $FF$ sont équiprobables. « Obtenir une pile et une face » contient $PF$ et $FP$ : sa probabilité vaut $\frac24=\frac12$, et non $\frac13$.\end{exemple}
\begin{definition}[Fréquences observées]Quand on répète une expérience, la fréquence d’une issue est son effectif divisé par le nombre de répétitions. Pour un même nombre de répétitions, les fréquences varient d’une série à l’autre : c’est la fluctuation. Une fréquence observée fournit une approximation de la probabilité.\end{definition}
\begin{attention}On distingue la situation réelle du modèle choisi pour la décrire : l’équiprobabilité se justifie par la symétrie d’une pièce ou d’un dé, pas par un résultat observé. Pour deux dés lancés ensemble, on les distingue, par exemple par leurs couleurs.\end{attention}

\clearpage\rubrique{Découverte}
\begin{blocLivre}{ex-11-1}\begin{exercice}[1]{Décrire un univers}
On lance un dé équilibré à huit faces numérotées de 1 à 8. Écrire l’univers, puis l’événement $A$ : « obtenir un multiple de 3 ». Calculer $P(A)$.
\end{exercice}\end{blocLivre}
\begin{blocLivre}{ex-11-2}\begin{exercice}[1]{L’événement contraire}
Avec le même dé à huit faces, écrire $\overline{A}$ et calculer sa probabilité de deux façons.
\end{exercice}\end{blocLivre}
\begin{blocLivre}{ex-11-3}\begin{exercice}[1]{Certain et impossible}
On lance un dé classique à six faces. Donner un événement certain, puis un événement égal à l’ensemble vide, avec leurs probabilités.
\end{exercice}\end{blocLivre}
\begin{blocLivre}{ex-11-4}\begin{exercice}[1]{Plusieurs écritures}
Une probabilité vaut $\frac38$. L’écrire sous forme décimale et en pourcentage.
\end{exercice}\end{blocLivre}
\begin{blocLivre}{ex-11-5}\begin{exercice}[1]{Une urne}
Une urne contient 4 boules rouges, 3 bleues et 1 verte, indiscernables au toucher. On tire une boule au hasard. On note $R$ l’événement « obtenir une boule rouge ». Calculer $P(R)$ et $P(\overline R)$.
\label{fig-ex-11-5}\figureGeo{urne-431}
\end{exercice}\end{blocLivre}

\rubrique{Entraînement}
\begin{blocLivre}{ex-11-6}\begin{exercice}[2]{Intersection et réunion}
On lance un dé équilibré à six faces. $A$ : « obtenir un nombre pair » ; $B$ : « obtenir au moins 4 ». Écrire $A$, $B$, $A\cap B$ et $A\cup B$, puis calculer leurs probabilités.
\label{fig-ex-11-6}\figureGeo{venn-de}
\end{exercice}\end{blocLivre}
\begin{blocLivre}{ex-11-7}\begin{exercice}[2]{Deux pièces}
On lance une pièce de 1 euro et une pièce de 2 euros, toutes deux équilibrées. Construire un arbre, donner les issues et calculer la probabilité d’obtenir une pile et une face.
\label{fig-ex-11-7}\figureGeo{arbre-deux-pieces}
\end{exercice}\end{blocLivre}
\begin{blocLivre}{ex-11-8}\begin{exercice}[2]{Une pièce et un dé}
On lance une pièce équilibrée puis un dé équilibré à six faces. Combien y a-t-il d’issues ? Calculer la probabilité de « pile et un nombre pair », puis celle de « face ou 6 ».
\label{fig-ex-11-8}\figureGeo{tableau-piece-de}
\end{exercice}\end{blocLivre}
\begin{blocLivre}{ex-11-9}\begin{exercice}[2]{La somme de deux dés}
On lance un dé rouge et un dé bleu, équilibrés, et l’on additionne les deux nombres. Compléter le tableau des sommes. Calculer la probabilité d’obtenir 7, puis celle d’obtenir au moins 10.
\label{fig-ex-11-9}\figureGeo{somme-deux-des}
\end{exercice}\end{blocLivre}
\begin{blocLivre}{ex-11-10}\begin{exercice}[2]{Des issues non équiprobables}
Un objet lancé a trois issues A, B, C de probabilités respectives 0,2 ; 0,5 et $p$. Trouver $p$, puis la probabilité de ne pas obtenir B.
\end{exercice}\end{blocLivre}

\rubrique{Problèmes}
\begin{blocLivre}{ex-11-11}\begin{exercice}[3]{Une observation}
Sur 50 lancers d’une pièce équilibrée, on observe 31 piles. Calculer la fréquence observée. Doit-on dire que la probabilité du prochain pile vaut $\frac{31}{50}$ ?
\end{exercice}\end{blocLivre}
\begin{blocLivre}{ex-11-12}\begin{exercice}[3]{Quatre séries de lancers}
Quatre groupes lancent chacun 50 fois une pièce équilibrée et obtiennent 23, 28, 25 et 31 piles. Calculer les quatre fréquences de pile. Pourquoi sont-elles différentes ? Calculer la fréquence sur les 200 lancers réunis.
\label{fig-ex-11-12}\figureGeo{fluctuation-pieces}
\end{exercice}\end{blocLivre}
\begin{blocLivre}{ex-11-13}\begin{exercice}[3]{Fréquences et probabilités}
On lance 60 fois deux pièces équilibrées et l’on compte les piles : 0 pile 13 fois, 1 pile 32 fois, 2 piles 15 fois. Calculer les fréquences au centième, puis compléter le diagramme en les comparant aux probabilités du modèle.
\label{fig-ex-11-13}\figureGeo{distribution-deux-pieces}
\end{exercice}\end{blocLivre}
\begin{blocLivre}{ex-11-14}\begin{exercice}[3]{Un effectif inconnu}
Un sac contient 5 boules rouges et un nombre $b$ de boules bleues, indiscernables au toucher. La probabilité de tirer une boule rouge vaut $\frac14$. Trouver $b$.
\end{exercice}\end{blocLivre}
\begin{blocLivre}{ex-11-15}\begin{exercice}[3]{Des numéros et des couleurs}
Un sac contient six jetons équiprobables : R1, R2, R3, B1, B2 et V1. $R$ : « le jeton est rouge » ; $U$ : « le jeton porte le numéro 1 ». Écrire $R\cap U$ et $R\cup U$, puis calculer $P(R\cup U)$. Pourquoi n’est-elle pas égale à $P(R)+P(U)$ ?
\label{fig-ex-11-15}\figureGeo{jetons-numerotes}
\end{exercice}\end{blocLivre}

\clearpage
